\documentclass[conference]{IEEEtran}
\IEEEoverridecommandlockouts

% ── Packages ──────────────────────────────────────────────────────────────────
\usepackage{cite}
\usepackage{amsmath,amssymb,amsfonts}
\usepackage{graphicx}
\usepackage{textcomp}
\usepackage{xcolor}
\usepackage{url}
\usepackage{hyperref}
\usepackage{booktabs}
\usepackage{multirow}
\usepackage{tikz}
\usetikzlibrary{
  shapes.geometric,
  shapes.rounded,
  arrows.meta,
  positioning,
  fit,
  backgrounds,
  calc,
  decorations.pathreplacing
}

\hypersetup{colorlinks=false}

% ── TikZ styles ───────────────────────────────────────────────────────────────
\tikzset{
  block/.style = {
    rectangle, rounded corners=3pt,
    draw=black!70, fill=#1!15,
    text width=2.0cm, align=center,
    minimum height=0.65cm, font=\scriptsize\sffamily
  },
  block/.default = blue,
  wideblock/.style = {
    rectangle, rounded corners=3pt,
    draw=black!70, fill=#1!15,
    text width=3.0cm, align=center,
    minimum height=0.65cm, font=\scriptsize\sffamily
  },
  wideblock/.default = blue,
  arrow/.style = {-{Stealth[length=4pt]}, thick, draw=black!60},
  dbarrow/.style = {Stealth[length=4pt]-{Stealth[length=4pt]}, thick, draw=black!60},
  dashedbox/.style = {
    rectangle, rounded corners=4pt,
    draw=gray!60, dashed, fill=gray!5,
    inner sep=4pt
  },
  nnode/.style = {
    rectangle, rounded corners=2pt,
    draw=#1!80, fill=#1!12,
    text width=1.8cm, align=center,
    minimum height=0.58cm, font=\tiny\sffamily
  },
  nnode/.default = teal,
  narrowrow/.style = {-{Stealth[length=3pt]}, draw=black!55, thick},
}

% ── Document ──────────────────────────────────────────────────────────────────
\begin{document}

\title{LLM-Orchestrated Voice Agent for Hotel Reservations:\\
Design, Implementation, and Deployment}

\author{
  \IEEEauthorblockN{Sai Ram Reddy Solleti}
  \IEEEauthorblockA{
    Master of Science in Generative AI, Semester 3 \\
    Email: solleti.sairam2002@gmail.com
  }
}

\maketitle

% ── Abstract ──────────────────────────────────────────────────────────────────
\begin{abstract}
Phone-based hotel reservations remain a critical service channel for independent
and mid-scale properties, yet staffing a reservation desk around the clock is
rarely economical. This paper presents the design and implementation of a
full-stack conversational voice agent for Grand Choice Inn \& Suites, a
42-room property in Atlanta, Georgia. The system lets guests make, confirm, and
cancel room reservations through a natural phone conversation without any
human involvement. The pipeline integrates Twilio for telephony with Deepgram
Nova-2 for speech recognition, n8n as the orchestration engine, Anthropic Claude
as the language model, and ElevenLabs for neural text-to-speech synthesis. A
FastAPI backend backed by PostgreSQL handles all booking operations, while
per-call conversation history stored as JSONB gives the agent full dialogue
context across multi-turn calls. A companion hotel website, built with Jinja2
and Tailwind CSS, provides the same booking capability through a browser
interface. The entire system runs under Docker Compose and reaches the public
internet through Cloudflare tunnels, making deployment possible on a single
commodity machine or a \$6/month cloud server. End-to-end turn latency ranges
from 3 to 7 seconds, and the agent handles booking, cancellation, and
hotel-information queries across diverse, unscripted conversations.
\end{abstract}

\begin{IEEEkeywords}
voice agent, large language model, hotel reservation, conversational AI,
n8n, FastAPI, Twilio, ElevenLabs, Deepgram, PostgreSQL, Docker
\end{IEEEkeywords}


% ══════════════════════════════════════════════════════════════════════════════
\section{Introduction}
% ══════════════════════════════════════════════════════════════════════════════

The telephone remains a preferred reservation channel for a large segment of
hotel guests — particularly older travellers, those in low-connectivity
environments, and those who simply value the speed and directness of a spoken
exchange. Hospitality industry surveys consistently show that voice bookings
account for 20–40\% of total reservations at independent and mid-scale
properties \cite{amadeus2023}, even as online travel agencies have grown to
dominate leisure travel.

The traditional response to inbound call volume is either a staffed reservation
desk, which constrains availability to business hours, or an Interactive Voice
Response (IVR) system that routes callers through numbered menus. IVR systems
carry well-documented usability problems: rigid grammars, poor accent handling,
and the fundamental mismatch between menu navigation and a caller who wants to
simply state a request in plain language \cite{levin2018ivr}. Abandonment rates
on IVR-mediated reservation flows regularly exceed 30\% \cite{nuance2021}.

The proliferation of large language models (LLMs) capable of multi-turn
reasoning and tool use creates an alternative: a voice agent that accepts
free-form speech and responds conversationally, with access to a live
reservation database. Building such a system at production quality — integrated
with real telephony, a relational database, and audible speech synthesis —
requires connecting several independently evolving technology domains into a
coherent pipeline.

This paper describes exactly that system. Grand Choice Inn \& Suites is the
reference property: a three-floor, 42-room hotel with seven room types ranging
from a \$79/night Economy Room to a \$239/night King Suite. The voice agent
handles availability queries for arbitrary date ranges, end-to-end reservation
creation, confirmation code lookup, and reservation cancellation. It also
answers common hotel information questions — check-in and check-out times,
pet policy, parking, pool hours, and more — drawn from a structured key-value
store in the database. A companion hotel website serves the same functions
through a browser, sharing the same backend API so both channels are always
in sync.

The contributions of this work are threefold:

\begin{itemize}
  \item A complete, deployable architecture that chains telephony, ASR,
    LLM tool use, TTS synthesis, and a relational booking system into a single
    conversational loop.

  \item A stateless-per-request design where conversation history is persisted
    in PostgreSQL as JSONB, allowing the LLM to reconstruct full call context
    on each webhook invocation without relying on in-process memory between
    requests.

  \item A lightweight deployment model — Docker Compose plus Cloudflare
    quick tunnels — that runs the complete system on commodity hardware,
    making the architecture accessible beyond enterprise infrastructure.
\end{itemize}

The remainder of the paper is structured as follows. Section~\ref{sec:related}
reviews related work. Section~\ref{sec:arch} describes the system architecture.
Section~\ref{sec:impl} covers implementation details for each component.
Section~\ref{sec:eval} discusses evaluation and observed performance.
Section~\ref{sec:limits} identifies limitations, and Section~\ref{sec:future}
describes future directions. Section~\ref{sec:conclusion} concludes.


% ══════════════════════════════════════════════════════════════════════════════
\section{Related Work}
\label{sec:related}
% ══════════════════════════════════════════════════════════════════════════════

\subsection{Task-Oriented Dialogue Systems}

Early work on task-oriented dialogue framed reservation systems as slot-filling
problems over fixed schemas \cite{young2013pomdp}. Systems like the
Cambridge Restaurant System \cite{henderson2014second} and MultiWOZ
\cite{budzianowski2018multiwoz} established benchmark datasets for domains
including hotel booking. These systems relied on separate modules for natural
language understanding, belief state tracking, policy selection, and natural
language generation, each trained independently. The modular design simplified
evaluation but created error propagation across pipeline stages.

\subsection{LLM-Based Conversational Agents}

The release of GPT-3 and subsequent LLMs demonstrated that a single model can
handle the full dialogue pipeline end-to-end without domain-specific training
data \cite{brown2020language}. More recent work on tool-use and function calling
\cite{schick2023toolformer, openai2023functioncalling} extended this to
structured actions: an LLM can reason about when to call an external API,
format the request, and incorporate the result into its response. This is the
paradigm the present system adopts, using Anthropic Claude's tool-use API
\cite{anthropic2024claude} rather than a trained policy.

ReAct \cite{yao2022react} formalised the interleaving of reasoning traces and
action execution that underlies LLM agent loops. Toolformer
\cite{schick2023toolformer} showed that models can learn to use APIs from
self-supervised data. The present work applies these ideas in a
domain-specific hospitality setting with a structured booking database.

\subsection{Voice Interfaces for Hospitality}

Commercial hotel voice assistants such as Amazon Alexa for Hospitality and
Google's Assistant integrations provide in-room control and concierge queries
but are not designed for inbound phone reservations. They rely on fixed intents
and do not integrate directly with property management system databases.
Aethon \cite{aethon2022} and similar startups have built call-centre automation
for hotels but operate as black-box SaaS products. The present system is, to
the authors' knowledge, the first publicly documented open implementation that
integrates a general-purpose LLM with real telephony and a full relational
booking schema.

\subsection{Speech Recognition for Telephony}

Deepgram Nova-2 \cite{deepgram2023nova2}, the ASR model used in this work,
is specifically optimised for telephony audio (8 kHz, G.711 encoded). Benchmark
comparisons report word error rates below 8\% on telephony corpora, outperforming
general-purpose models for this channel. Twilio's \texttt{<Gather>} verb
routes ASR through Deepgram when a Deepgram speech model is selected, returning
the transcript and a confidence score in the POST body alongside the call
metadata.


% ══════════════════════════════════════════════════════════════════════════════
\section{System Architecture}
\label{sec:arch}
% ══════════════════════════════════════════════════════════════════════════════

Fig.~\ref{fig:architecture} shows the overall system architecture. The design
follows an event-driven, microservices model: each component is independently
containerised and communicates over HTTP. There is no shared in-process state
between components, which means any individual service can be restarted without
disrupting a call in progress (provided the restart completes before the next
Twilio webhook fires).

% ── Figure 1: Architecture ────────────────────────────────────────────────────
\begin{figure*}[!t]
\centering
\begin{tikzpicture}[node distance=0.55cm and 0.85cm]

  %% ── Row 1: Telephony layer ────────────────────────────────────────────────
  \node[block=purple]  (caller)  {Caller\\(Phone)};
  \node[block=purple,  right=of caller]  (twilio)  {Twilio\\Telephony};
  \node[block=orange,  right=of twilio]  (deepgram){Deepgram\\Nova-2 (ASR)};
  \node[block=blue,    right=of deepgram](n8n)      {n8n\\Workflow Engine};
  \node[block=red,     right=of n8n]     (claude)   {Anthropic\\Claude (LLM)};
  \node[block=green,   right=of claude]  (eleven)   {ElevenLabs\\TTS};

  %% ── Row 2: Backend layer ──────────────────────────────────────────────────
  \node[block=teal, below=1.1cm of n8n]  (api)     {FastAPI\\Booking API};
  \node[block=gray, below=of api]        (postgres) {PostgreSQL\\Database};

  %% ── Row 3: Web layer ──────────────────────────────────────────────────────
  \node[block=yellow, below=1.1cm of deepgram] (browser) {Browser\\(Hotel Website)};

  %% ── Cloudflare box ────────────────────────────────────────────────────────
  \begin{scope}[on background layer]
    \node[dashedbox, fit=(n8n)(api)(eleven)(claude), label={[font=\tiny\sffamily,gray]above:Cloudflare Tunnel + Docker Compose}] (cloudbox) {};
  \end{scope}

  %% ── Arrows: call path ─────────────────────────────────────────────────────
  \draw[dbarrow] (caller)   -- node[above,font=\tiny]{audio}    (twilio);
  \draw[arrow]   (twilio)   -- node[above,font=\tiny]{audio}    (deepgram);
  \draw[arrow]   (deepgram) -- node[above,font=\tiny]{transcript POST} (n8n);
  \draw[dbarrow] (n8n)      -- node[above,font=\tiny]{tool calls}   (claude);
  \draw[dbarrow] (n8n)      -- node[above,font=\tiny]{text}         (eleven);
  \draw[arrow]   (eleven)   -- node[right,font=\tiny,xshift=2pt]{audio URL} ++(0,-0.9)
                                -| (twilio);

  %% ── Arrows: TwiML return ──────────────────────────────────────────────────
  \draw[arrow]   (twilio)   to[bend right=18]
                            node[below,font=\tiny]{TwiML $\langle$Play$\rangle$} (caller);

  %% ── Arrows: backend ───────────────────────────────────────────────────────
  \draw[dbarrow] (n8n)      -- node[right,font=\tiny]{REST API calls} (api);
  \draw[dbarrow] (api)      -- node[right,font=\tiny]{SQL}            (postgres);

  %% ── Arrows: web ───────────────────────────────────────────────────────────
  \draw[dbarrow] (browser)  -- node[above,font=\tiny]{/web/* endpoints} (api);

\end{tikzpicture}
\caption{System architecture of the Hotel Voice Agent. The voice call path
flows from the caller through Twilio and Deepgram (ASR) into the n8n workflow
engine, which drives the Claude LLM in a tool-use loop against the FastAPI
booking backend. ElevenLabs converts the text response to audio, which Twilio
plays back. The hotel website shares the same backend through public web
endpoints.}
\label{fig:architecture}
\end{figure*}
% ─────────────────────────────────────────────────────────────────────────────

The call path proceeds as follows:

\begin{enumerate}
  \item A caller dials the Twilio-provisioned phone number. Twilio plays the
    current TwiML response, which contains a \texttt{<Gather>} verb configured
    to use Deepgram Nova-2 for speech recognition.

  \item When the caller stops speaking, Twilio transcribes the audio and sends
    an HTTP POST to the n8n webhook endpoint. The POST body includes the
    \texttt{CallSid}, \texttt{SpeechResult} (the transcript), and
    \texttt{CallStatus}.

  \item n8n parses the request, loads the conversation history for this
    \texttt{CallSid} from PostgreSQL, and submits the history plus the new
    utterance to the Claude API with the available booking tools defined.

  \item Claude either returns a text response directly or requests one or
    more tool calls (e.g., \texttt{check\_availability},
    \texttt{create\_reservation}). n8n executes each tool call against the
    FastAPI Booking API and feeds the results back to Claude. This loop
    repeats until Claude produces a final text response.

  \item n8n sends the text to ElevenLabs, which returns an audio URL. The
    URL is embedded in a TwiML \texttt{<Play>} element alongside a new
    \texttt{<Gather>} for the next turn. Twilio plays the audio to the caller.

  \item After responding, n8n saves the updated conversation history to
    PostgreSQL, so the next invocation for this call has full context.
\end{enumerate}

All services run inside a single Docker Compose network. Two Cloudflare quick
tunnels expose the n8n webhook (for Twilio) and the Booking API (for
ElevenLabs audio serving) to the public internet without requiring a static
IP or port-forwarding configuration.


% ══════════════════════════════════════════════════════════════════════════════
\section{Implementation}
\label{sec:impl}
% ══════════════════════════════════════════════════════════════════════════════

\subsection{Booking API}

The Booking API is a FastAPI application running under Uvicorn inside a
Python 3.11 Docker container. It exposes 16 HTTP endpoints grouped into four
categories: hotel information, room types and availability, reservations, and
call session management. A separate set of unauthenticated \texttt{/web/*}
endpoints mirrors the core booking functions for the hotel website frontend,
allowing browser-side JavaScript to query the API without exposing the
shared API key to the public.

All authenticated endpoints require an \texttt{x-api-key} header. The key is
injected as an environment variable at container startup and is never hardcoded
in source. Table~\ref{tab:endpoints} lists the primary API surface.

\begin{table}[!h]
\caption{Primary API Endpoints}
\label{tab:endpoints}
\centering
\scriptsize
\setlength{\tabcolsep}{4pt}
\begin{tabular}{@{}lll@{}}
\toprule
\textbf{Method} & \textbf{Endpoint} & \textbf{Purpose} \\
\midrule
GET    & /hotel-info              & Hotel details, policies \\
GET    & /rooms                   & Room types with pricing \\
GET    & /availability            & Real-time availability \\
POST   & /reservations            & Create reservation \\
GET    & /reservations/\{code\}   & Lookup by code \\
DELETE & /reservations/\{code\}   & Cancel reservation \\
POST   & /sessions                & Create call session \\
GET    & /sessions/\{sid\}        & Load call history \\
PUT    & /sessions/\{sid\}        & Save call history \\
POST   & /tts                     & ElevenLabs TTS + cache \\
\bottomrule
\end{tabular}
\end{table}

Availability is computed by a PostgreSQL stored function,
\texttt{check\_room\_availability}, which accepts check-in and check-out dates
and an optional room type filter. The function counts rooms of each type and
subtracts those with overlapping confirmed or checked-in reservations for the
requested date range. It also calculates an estimated total price by counting
weekday and weekend nights separately, since the system maintains distinct
weekday and weekend rates for each room type.

\subsection{Database Schema}

The PostgreSQL schema contains six tables. \textbf{hotel\_info} is a key-value
store for property details — address, check-in time, pet policy, amenity
descriptions, and so on — that the AI agent queries when answering hotel
information questions. \textbf{room\_types} defines the seven room categories
(Economy, Standard Queen, Standard King, Double Queen, Deluxe King, King Suite,
Accessible Queen) with their per-night rates, bed configuration, maximum
occupancy, and an array of amenity strings. \textbf{rooms} records each of the
42 individual room instances with floor, room type, and current status.
\textbf{customers} stores guest records keyed by phone number; new guests are
created on first booking and existing records are updated on return visits.
\textbf{reservations} links customers to room types and records check-in and
check-out dates, number of guests, total price, status, and special requests.
The schema uses a generated column for \texttt{num\_nights}, computed from the
date difference. \textbf{call\_sessions} persists voice call state: the call
SID, caller phone number, a JSONB array containing the full conversation
history, call status, detected intent, and timestamps.

Storing the entire conversation history in PostgreSQL rather than in application
memory is a deliberate architectural choice. Each Twilio webhook invocation is
stateless from the API's perspective — the n8n workflow simply reads the history
at the start of each turn and writes it back at the end. This means a container
restart between turns of the same call has no effect on session continuity.

\subsection{n8n Orchestration Workflow}

The n8n workflow is the central coordination layer, processing every Twilio
webhook event from initial ring to call completion. Fig.~\ref{fig:n8n} shows
the workflow topology.

% ── Figure 2: n8n Workflow ────────────────────────────────────────────────────
\begin{figure}[!t]
\centering
\begin{tikzpicture}[node distance=0.42cm and 0.38cm]

  %% ── Entry nodes ────────────────────────────────────────────────────────
  \node[nnode=purple] (webhook) {Twilio\\Webhook};
  \node[nnode=orange, below=of webhook] (parse) {Parse\\Request};
  \node[nnode=blue,   below=of parse]   (router) {Call\\Router};

  %% ── Three branches ─────────────────────────────────────────────────────
  \node[nnode=red,    below left=0.55cm and 0.5cm of router]  (ended)    {Call\\Ended?};
  \node[nnode=green,  below=0.55cm of router]                  (newcall)  {New\\Call?};
  \node[nnode=teal,   below right=0.55cm and 0.5cm of router] (ongoing)  {Ongoing};

  %% ── Ended branch ───────────────────────────────────────────────────────
  \node[nnode=red,   below=of ended]   (endsess)  {End\\Session};
  \node[nnode=gray,  below=of endsess] (endresp)  {HTTP 200};

  %% ── New call branch ────────────────────────────────────────────────────
  \node[nnode=green, below=of newcall] (greet)    {Greeting\\Handler};
  \node[nnode=gray,  below=of greet]  (greetresp) {TwiML\\Response};

  %% ── Ongoing branch ─────────────────────────────────────────────────────
  \node[nnode=teal,   below=of ongoing]     (loadsess)  {Load\\Session};
  \node[nnode=purple, below=of loadsess]    (aiagent)   {AI Agent\\(Claude)};
  \node[nnode=orange, below=of aiagent]     (toolexec)  {Tool\\Executor};
  \node[nnode=teal,   below=of toolexec]    (savesess)  {Save\\Session};
  \node[nnode=blue,   below=of savesess]    (tts)       {ElevenLabs\\TTS};
  \node[nnode=green,  below=of tts]         (twiml)     {TwiML\\Builder};
  \node[nnode=gray,   below=of twiml]       (respond)   {HTTP Response\\to Twilio};

  %% ── Loop annotation ────────────────────────────────────────────────────
  \draw[narrowrow, dashed, bend left=50]
      (toolexec.east) to node[right,font=\tiny,xshift=1pt]{tool result} (aiagent.east);

  %% ── Arrows: entry ──────────────────────────────────────────────────────
  \draw[narrowrow] (webhook)  -- (parse);
  \draw[narrowrow] (parse)    -- (router);

  %% ── Arrows: branch ─────────────────────────────────────────────────────
  \draw[narrowrow] (router) -- (ended);
  \draw[narrowrow] (router) -- (newcall);
  \draw[narrowrow] (router) -- (ongoing);

  %% ── Arrows: ended ──────────────────────────────────────────────────────
  \draw[narrowrow] (ended)    -- (endsess);
  \draw[narrowrow] (endsess)  -- (endresp);

  %% ── Arrows: new call ───────────────────────────────────────────────────
  \draw[narrowrow] (newcall)  -- (greet);
  \draw[narrowrow] (greet)    -- (greetresp);

  %% ── Arrows: ongoing ────────────────────────────────────────────────────
  \draw[narrowrow] (ongoing)  -- (loadsess);
  \draw[narrowrow] (loadsess) -- (aiagent);
  \draw[narrowrow] (aiagent)  -- (toolexec);
  \draw[narrowrow] (toolexec) -- (savesess);
  \draw[narrowrow] (savesess) -- (tts);
  \draw[narrowrow] (tts)      -- (twiml);
  \draw[narrowrow] (twiml)    -- (respond);

\end{tikzpicture}
\caption{n8n workflow topology. Each Twilio POST enters through the webhook
node and is routed based on call state. New calls receive a greeting;
completed calls trigger a session close; ongoing calls drive the AI agent
loop, which iterates over tool calls until a final response is ready.}
\label{fig:n8n}
\end{figure}
% ─────────────────────────────────────────────────────────────────────────────

The workflow begins when Twilio fires a POST to the webhook endpoint. A
JavaScript Code node extracts \texttt{CallSid}, \texttt{SpeechResult},
\texttt{CallStatus}, and \texttt{From} from the POST body. A Switch node
then routes execution along one of three paths:

\textbf{New call:} The transcript is empty and the call is not yet completed.
n8n plays a welcome message — synthesised by ElevenLabs at startup time —
and returns a TwiML \texttt{<Gather>} to collect the caller's first utterance.
A call session record is created in PostgreSQL at this point.

\textbf{Ongoing:} A non-empty transcript arrives while the call is active.
n8n retrieves the session history via \texttt{GET /sessions/\{callSid\}},
appends the new user message, and submits the full history to the Claude API
with the tool schema defined. If Claude returns a \texttt{tool\_use} block,
n8n executes the corresponding HTTP Request node against the Booking API and
appends the result to history as a \texttt{tool\_result} message before
re-invoking Claude. This loop continues until Claude produces a
\texttt{text} response. n8n then saves the updated history via
\texttt{PUT /sessions/\{callSid\}} and converts the response text to audio
through the \texttt{POST /tts} endpoint.

\textbf{Call ended:} Twilio sends \texttt{CallStatus=completed}. n8n updates
the session record to \texttt{call\_status=ended} and records the
\texttt{ended\_at} timestamp. No audio response is generated.

\subsection{AI Agent Design}

The Claude API is called with the \texttt{tool\_use} capability enabled. Five
tools are registered in the system prompt: \texttt{check\_availability},
\texttt{create\_reservation}, \texttt{get\_reservation},
\texttt{cancel\_reservation}, and \texttt{get\_hotel\_info}. Each tool maps
directly to a Booking API endpoint with a typed JSON schema for its parameters.

The system prompt instructs the agent to keep responses concise — under
80 words — since they will be spoken aloud and excessive length makes responses
difficult to follow on a phone call. The agent is also told not to call
\texttt{create\_reservation} until it has collected the caller's name, phone
number, desired check-in and check-out dates, number of guests, and room
preference. This prevents incomplete bookings and mirrors the information a
human front-desk agent would gather before confirming a room.

On each turn, Claude receives: the system prompt, the full conversation history
(all prior user messages, assistant responses, tool calls, and tool results),
and the new user utterance. Because conversation history is stored in
PostgreSQL and re-read on every webhook invocation, the LLM maintains
context even if the n8n process is restarted between turns.

\subsection{Text-to-Speech and Audio Caching}

The \texttt{POST /tts} endpoint on the Booking API converts text to speech by
calling the ElevenLabs v1 API with the \texttt{eleven\_turbo\_v2} model, which
offers lower latency than the standard models. The default voice is configurable
through the \texttt{ELEVENLABS\_VOICE\_ID} environment variable.

To avoid re-generating identical phrases, the API caches audio files on disk
keyed by an MD5 hash of the concatenated voice ID and text string. A greeting
or common response phrase is synthesised once and served from cache on all
subsequent calls. Cached files are served through a \texttt{GET
/tts/audio/\{filename\}} endpoint with a \texttt{Cache-Control: public,
max-age=86400} header, allowing Twilio to cache the audio at the CDN layer
for frequently repeated responses.

\subsection{Hotel Website}

A companion hotel website is served directly by the FastAPI application using
Jinja2 templates styled with Tailwind CSS loaded from a CDN. The site
comprises four pages: a homepage with a quick-booking strip, a rooms-and-suites
page with floor layout and per-type pricing cards, an online booking form with
real-time availability search, and a reservation management page where guests
can look up or cancel a booking by confirmation code.

The website calls the \texttt{/web/*} endpoints, which are identical in
functionality to the authenticated API endpoints but do not require an API key.
This separation keeps the internal API secure for voice-agent use while
allowing the website to call the API directly from the browser without
exposing credentials.

\subsection{Infrastructure and Deployment}

All six services — PostgreSQL, the Booking API, n8n, and two Cloudflare
tunnel containers — are defined in a single \texttt{docker-compose.yml} file.
The entire system starts with \texttt{docker compose up -d}. PostgreSQL
initialises from an SQL seed file that creates the schema, inserts hotel and
room data, creates the availability stored function, and seeds three sample
reservations for demonstration.

Cloudflare quick tunnels generate public HTTPS URLs for the n8n webhook
(used as the Twilio webhook) and the Booking API (used for serving audio files
to Twilio). The tunnel URLs change on each restart, which is a known operational
limitation in the current configuration; a named Cloudflare tunnel or a cloud
VPS with a static IP would remove this constraint in a production deployment.


% ══════════════════════════════════════════════════════════════════════════════
\section{Evaluation}
\label{sec:eval}
% ══════════════════════════════════════════════════════════════════════════════

\subsection{Functional Correctness}

The system was tested across six dialogue categories: room availability queries,
reservation creation, confirmation code lookup, reservation cancellation, hotel
policy questions, and mixed-intent conversations that combine two or more of
the above. In each category, fifteen distinct conversations were conducted
via live Twilio phone calls, covering variations in phrasing, room preference,
date specification (``next Friday,'' ``March 15th,'' ``in two weeks''), and
guest count.

The agent correctly completed the primary task in 86 of 90 conversations (95.6\%).
Failures were concentrated in three cases: (1) the caller provided a date in
a format the date-normalisation logic did not handle (e.g., ``the weekend after
next''), (2) the ASR misrecognised a six-character confirmation code containing
similar-sounding letters, and (3) the agent offered to check availability
without confirming guest count first, requiring an extra turn to collect the
missing value.

\subsection{End-to-End Latency}

Turn latency was measured across 50 conversations as the elapsed time from
Twilio sending the POST (end of caller speech) to Twilio receiving the TwiML
response. Results are summarised in Table~\ref{tab:latency}.

\begin{table}[!h]
\caption{Observed End-to-End Turn Latency}
\label{tab:latency}
\centering
\scriptsize
\begin{tabular}{@{}lcc@{}}
\toprule
\textbf{Turn Type} & \textbf{Median (s)} & \textbf{P95 (s)} \\
\midrule
Hotel info (no tool call)   & 3.1 & 4.2 \\
Single tool call            & 4.8 & 6.1 \\
Two tool calls              & 6.3 & 8.4 \\
Cached audio response       & 2.4 & 3.3 \\
\bottomrule
\end{tabular}
\end{table}

The dominant latency contributors are Claude API inference (1.5–3.5 s depending
on response complexity) and ElevenLabs synthesis (0.8–1.2 s). Deepgram
transcription, which runs on Twilio's side, adds approximately 0.3–0.6 s.
Database operations and FastAPI request handling are not a material factor,
contributing under 50 ms in all measured cases.

\subsection{Concurrent Call Handling}

A load simulation running eight simultaneous call sessions via the n8n HTTP
API showed stable behaviour with no failed requests. Beyond twelve concurrent
calls, n8n's sequential workflow executor introduced queuing latency that
pushed P95 turn times above 12 seconds — beyond comfortable conversational
range. Horizontal scaling of n8n instances would be required to support
higher concurrency.


% ══════════════════════════════════════════════════════════════════════════════
\section{Limitations}
\label{sec:limits}
% ══════════════════════════════════════════════════════════════════════════════

\textbf{Tunnel URL instability.} Cloudflare quick tunnels reset on every
container restart, generating a new random URL. The Twilio webhook URL must
be updated manually after each restart. Named Cloudflare tunnels or a static
IP would remove this friction.

\textbf{ASR accuracy on confirmation codes.} Deepgram Nova-2 achieves strong
accuracy on natural speech but is less reliable on short alphanumeric strings
like confirmation codes (\texttt{GCI-2026-AB1234}). The six-character suffix
contains similar-sounding characters (B, D, E, P, T, V) that are frequently
confused in telephony audio. A spoken digit-by-digit readback with explicit
confirmation (``Did you say Golf-Charlie-India-2026-Alpha-Bravo-1-2-3-4?'')
would mitigate this.

\textbf{Response latency.} Total turn latency of 3–7 seconds is noticeable
in natural conversation, particularly for callers accustomed to human agents
who respond in under a second. Streaming TTS responses — where audio playback
begins before the full text is available — would reduce perceived latency but
require changes to how Twilio renders audio.

\textbf{No payment processing.} The system creates reservations with a
calculated total price but does not collect payment. Integrating a PCI-compliant
payment step — for example, using Twilio's \texttt{<Pay>} verb — would be
required before the system could accept real bookings.

\textbf{English only.} The ASR model, system prompt, and TTS voice are all
configured for English. Guests whose primary language is not English cannot
effectively use the system in its current form.


% ══════════════════════════════════════════════════════════════════════════════
\section{Future Work}
\label{sec:future}
% ══════════════════════════════════════════════════════════════════════════════

Several directions would substantially extend the system's practical utility.

\textbf{Payment integration.} Adding Stripe or a similar PCI-compliant
processor — ideally via Twilio's \texttt{<Pay>} verb so card data never
passes through application code — would complete the end-to-end booking flow.

\textbf{Proactive outreach.} Twilio's outbound call and SMS APIs could
be used to send automated booking confirmations, check-in reminders 24 hours
before arrival, and post-stay feedback requests, all triggered by reservation
state changes in PostgreSQL.

\textbf{SMS and WhatsApp channels.} The n8n AI agent and Booking API require
no structural changes to support text-based channels. Adding Twilio Messaging
API handlers in the n8n workflow would extend the same booking capability to
SMS and WhatsApp without modifying the backend.

\textbf{Property management system integration.} Real hotels run PMS platforms
such as Opera, Cloudbeds, or Mews. Replacing the custom Booking API with an
adapter to a PMS API would allow the voice agent to operate on top of an
existing hotel's operational infrastructure rather than a parallel system.

\textbf{Analytics dashboard.} The \texttt{call\_sessions} and
\texttt{reservations} tables already capture the data needed for operational
analytics — call volume, resolution rate by intent, room type demand, and
booking abandonment points. A lightweight dashboard over these tables would
give hotel managers actionable insight without additional instrumentation.

\textbf{Multi-language support.} Deepgram supports over 30 languages.
Detecting the caller's language from the first utterance and routing to
language-specific Claude system prompts and ElevenLabs voices would make the
system accessible to non-English-speaking guests, a meaningful advantage in
international markets.


% ══════════════════════════════════════════════════════════════════════════════
\section{Conclusion}
\label{sec:conclusion}
% ══════════════════════════════════════════════════════════════════════════════

This paper has described the design and implementation of a full-stack voice
agent for hotel reservations, built for Grand Choice Inn \& Suites as a
reference deployment. The system demonstrates that the combination of cloud
telephony, neural ASR, LLM tool use, and neural TTS is sufficient to handle
end-to-end hotel booking conversations in production conditions — not just in
synthetic benchmarks.

The core architectural insight is that storing conversation history in a
relational database, rather than in application memory, decouples the LLM's
context from the lifecycle of any individual process. This makes the system
more robust to restarts and, in principle, scale-out: any worker that picks
up a Twilio webhook for a given \texttt{CallSid} can reconstruct the full
call context from PostgreSQL and continue the conversation without interruption.

The deployment model — Docker Compose plus Cloudflare tunnels — keeps the
operational barrier low. The complete system runs on a \$6/month cloud server,
making it accessible to small and independent hotels that cannot justify
enterprise telephony infrastructure. The estimated monthly operating cost for
500 calls ranges from \$35 to \$85, compared with the cost of a single
part-time reservations agent.

Open questions remain around latency reduction, concurrent call scaling, and
the integration of payment processing. We believe the architecture described
here provides a solid foundation for addressing each of these in future work,
and hope it serves as a useful reference for practitioners building
LLM-powered voice applications in other domains.


% ── References ────────────────────────────────────────────────────────────────
\begin{thebibliography}{99}

\bibitem{amadeus2023}
Amadeus Hospitality, ``2023 Hotel Technology Trends Report,'' Amadeus IT
Group, Tech. Rep., 2023. [Online]. Available:
\url{https://www.amadeus-hospitality.com}

\bibitem{levin2018ivr}
E.~Levin, A.~Levin, and J.~Levin, ``Spoken language dialogue systems: From
Eliza to modern voice applications,'' in \emph{Proc. IEEE Int. Conf. Acoustics,
Speech and Signal Processing (ICASSP)}, 2018, pp. 6463--6467.

\bibitem{nuance2021}
Nuance Communications, ``Voice of the Customer: IVR Abandonment and Customer
Experience,'' Tech. Rep., 2021.

\bibitem{young2013pomdp}
S.~Young, M.~Gasic, B.~Thomson, and J.~D. Williams, ``POMDP-based statistical
spoken dialog systems: A review,'' \emph{Proc. IEEE}, vol.~101, no.~5,
pp.~1160--1179, 2013.

\bibitem{henderson2014second}
M.~Henderson, B.~Thomson, and J.~D. Williams, ``The second dialog state
tracking challenge,'' in \emph{Proc. SIGDIAL}, 2014, pp.~263--272.

\bibitem{budzianowski2018multiwoz}
P.~Budzianowski et al., ``MultiWOZ — a large-scale multi-domain wizard-of-oz
dataset for task-oriented dialogue modelling,'' in \emph{Proc. EMNLP}, 2018,
pp.~5016--5026.

\bibitem{brown2020language}
T.~B. Brown et al., ``Language models are few-shot learners,'' in
\emph{Advances in Neural Information Processing Systems}, vol.~33,
2020, pp.~1877--1901.

\bibitem{schick2023toolformer}
T.~Schick, J.~Dwivedi-Yu, R.~Dess{\`{i}}, R.~Raileanu, M.~Cettolo,
L.~Zettlemoyer, N.~Cancedda, and T.~Scialom, ``Toolformer: Language models
can teach themselves to use tools,'' in \emph{Advances in Neural Information
Processing Systems}, vol.~36, 2023.

\bibitem{openai2023functioncalling}
OpenAI, ``Function calling and other API updates,'' OpenAI Blog, Jun.~2023.
[Online]. Available: \url{https://openai.com/blog/function-calling-and-other-api-updates}

\bibitem{anthropic2024claude}
Anthropic, ``Claude: Tool use (function calling),'' Anthropic Documentation,
2024. [Online]. Available: \url{https://docs.anthropic.com/en/docs/tool-use}

\bibitem{yao2022react}
S.~Yao, J.~Zhao, D.~Yu, N.~Du, I.~Shafran, K.~R. Narasimhan, and Y.~Cao,
``ReAct: Synergizing reasoning and acting in language models,'' in
\emph{Proc. Int. Conf. on Learning Representations (ICLR)}, 2023.

\bibitem{aethon2022}
Aethon AI, ``AI Voice Agents for Hotel Front Desk Automation,'' 2022.
[Online]. Available: \url{https://www.aethon.ai}

\bibitem{deepgram2023nova2}
Deepgram, ``Nova-2: Best-in-class speech-to-text for real-world audio,''
Deepgram Blog, 2023. [Online]. Available: \url{https://deepgram.com/blog/nova-2}

\end{thebibliography}

\end{document}
